\documentclass[openany]{book}

\usepackage[margin=1in]{geometry}
\usepackage{amsmath,amsfonts,amsthm, amssymb}
\usepackage{yhmath}
\usepackage{mathrsfs}
\usepackage{mathtools}
\usepackage{xcolor}
\usepackage{graphicx}
\usepackage{comment}
\usepackage{tikz-cd}
\usepackage{quiver}
\usepackage{hyperref,cleveref}
\renewcommand{\familydefault}{ppl}
\newcommand{\curl}{\text{curl}}
\newcommand{\diverg}{\text{div}}
\newcommand{\tr}{\text{tr}}
\newcommand{\R}{\mathbb{R}}
\newcommand{\E}{\mathbb{E}}
\newcommand{\Z}{\mathbb{Z}}
\newcommand{\C}{\mathbb{C}}
\newcommand{\F}{\mathbb{F}}
\newcommand{\la}{\langle}
\newcommand{\ra}{\rangle}
\newcommand{\colim}{\text{colim}}
\DeclareMathOperator{\im}{im}
\DeclareMathOperator{\disc}{disc}
\let\oldemptyset\emptyset
\let\emptyset\varnothing
\newcommand{\tor}{\text{Tor}}
\newcommand{\id}{\text{id}}
\newcommand{\ext}{\text{Ext}}
\newcommand{\ptop}{\text{PTop}}
\newcommand{\pt}{\text{pt}}
\newcommand{\ach}{\text{Ach}}
\newcommand{\Q}{\mathbb{Q}}
\newcommand{\gal}{\text{Gal}}
\newcommand{\fraccomma}{\genfrac{}{}{0pt}{}{}{,}}
\definecolor{wikipediadarkblue}{rgb}{0.023, 0.270, 0.676}
\hypersetup{
    colorlinks,
    citecolor=black,
    filecolor=black,
    linkcolor=blue,
    urlcolor=red
}

% \hypersetup{
%     urlbordercolor={1 0 0}, % Red box for URLs
%     linkbordercolor=blue, % Blue box for internal links
%     % the values are in RGB format from 0 to 1
% }


\input{hui_r.tex}

\title{Complex Analysis Qualifying Exam Review
\\ 
\vspace{0.4cm}
\large Fall 2026}



% \author{(Please email hsun95@jh.edu if you see typos)}
% \date{\today}


\begin{document}

\maketitle

\tableofcontents
\newpage



\chapter{Fall 2025}
\begin{prob}{5}
Let $\{f_n\}$ be a sequence of holomorphic functions on the open unit disk
\[
\mathbb{D}=\{z\in\mathbb{C}:|z|<1\},
\]
and suppose that $f_n\to f$ uniformly on compact subsets of $\mathbb{D}$, where
$f:\mathbb{D}\to\mathbb{C}$ is continuous. Prove that $f$ is holomorphic on $\mathbb{D}$.
\end{prob}
\begin{proof}
    Morera's theorem.
\end{proof}

\begin{prob}{6}
Use contour integration to evaluate
\[
\int_0^\infty \frac{x^2}{x^4+1}\,dx.
\]
\end{prob}
\begin{proof}
    Use residue theorem for $z_1=e^\frac{\pi i}{4}, z_2=e^\frac{3\pi i}{4}$.
\end{proof}

\begin{prob}{7}
Let $f$ be an analytic function on the open unit disk
\[
\mathbb{D}=\{z\in\mathbb{C}:|z|<1\},
\]
and suppose that
\[
|f(z)|\geq \sqrt{|z|}
\]
for all $z\in\mathbb{D}$. Show that
\[
|f(z)|\geq 1
\qquad\text{for all }z\in\mathbb{D}.
\]
\end{prob}
\begin{proof}
    Maximum modulus principle.
\end{proof}

\begin{prob}{8}
Determine how many solutions
\[
z^7-7z^4+z^3-3z=1
\]
are in the annulus
\[
1<|z|<2.
\]
\end{prob}
\begin{proof}
    Rouche's theorem.
\end{proof}

% \chapter{Spring 2025}
% \begin{prob}
% Let $f(z)$ be holomorphic on $\mathbb{C}\setminus\mathbb{R}$. $f$ is continuous at each point of $\mathbb{R}$. Show that $f$ extends to an entire analytic function. (Hint: Use Morera's theorem.)
% \end{prob}


% \begin{prob}
% Evaluate
% \[
% \operatorname{p.v.}\int_{-\infty}^{\infty}\frac{e^{ix}}{x^2}\,dx
% \]
% by residue theory.
% \end{prob}

% \begin{prob}
% Let $f:\Omega\to\mathbb{C}$ be a non-constant holomorphic function defined on a connected open set $\Omega\subset\mathbb{C}$, and suppose that $f$ is open (i.e., maps open sets to open sets). Additionally, suppose that for some $z_0\in\Omega$, we have $f'(z_0)=0$. Show that $f$ is not injective on $\Omega$.
% \end{prob}

% \begin{prob}
% Let $f:\mathbb{D}\to\mathbb{C}$ be a holomorphic function on the open unit disk
% \[
% \mathbb{D}=\{z\in\mathbb{C}:|z|<1\},
% \]
% and suppose that
% \[
% |f(z)|\leq M \qquad \text{for all }z\in\mathbb{D},
% \]
% for some constant $M>0$. Prove the estimate that for all $z\in\mathbb{D}$,
% \[
% |f'(z)|\leq\frac{M}{(1-|z|)^2}.
% \]
% \end{prob}



\chapter{Cauchy's integral formula}

\begin{prob}
    Evaluate the following integral:
    \begin{equation*}
        \int_0^\infty\frac{1-\cos x}{x^2}dx
    \end{equation*}
\end{prob}
% \begin{proof}
%     The answer is 
%     \begin{equation*}
%         \int_0^\infty\frac{1-\cos x}{x^2}dx=\frac{\pi}{2}
%     \end{equation*}
%     We use the contour of a semicircle skipping $x=0$. Let $f(z)=\frac{1-e^{iz}}{z^2}$, then by Cauchy's integral formula, we have 
%     \begin{align*}
%         \int_{\gamma_R}f(z)dz+\int_{\gamma_\varepsilon}f(z)dz+\int_{|z|>\varepsilon}f(z)dz=0\tag{1}
%     \end{align*}
% where 
% \begin{equation*}
%     \lim_{\varepsilon\to 0}\int_{|z|>\varepsilon}f(z)dz=2\int_0^\infty\frac{1-\cos x}{x^2}dx
% \end{equation*}
% Hence it suffices to compute the first two terms in (1). Let $z=Re^{i\theta}$, we see that 
% \begin{align*}
%     \left|\int_{\gamma_R}f(z)dz\right|&=\left|\int_0^{\pi}\frac{1-e^{iRe^{i\theta}}}{R^2e^{2i\theta}}iRe^{i\theta}d\theta\right|\\
%     &\leq\int_0^\pi\frac{2}{R}d\theta\to 0
% \end{align*}
% as $R\to\infty$.
% Similarly, we can compute $\int_{\gamma_\varepsilon}f(z)dz$ by setting $z=\varepsilon e^{i\theta}$ and letting $\varepsilon\to 0$. 
% \begin{align*}
%     \int_{\gamma_\varepsilon}f(z)dz&=\int_\pi^{0}\frac{1-e^{i\varepsilon e^{i\theta}}}{\varepsilon^2e^{2i\theta}}i\varepsilon e^{i\theta}d\theta\\
%     &=i\int_\pi^0\frac{1-e^{i\varepsilon e^{i\theta}}}{\varepsilon e^{i\theta}}d\theta\to -\pi
% \end{align*}
% as $\varepsilon\to 0$, since
% \begin{equation*}
%     \lim_{a\to 0}\frac{1-e^{ia}}{a}=-i
% \end{equation*}
% \end{proof}


\begin{prob}
    Let $U$ be a region. Suppose $f$ is holomorphic in $U$, and a sequence $\{z_n\}\subset U$ converges to some $z\in U$. If $f(z_n)=0$ for all $n$, show that $f\equiv 0$.
\end{prob}
% \begin{proof}
%     Since $z_n\to z$, we know $f(z)=\lim_{z_n\to z}f(z_n)=0$. Assume $f$ is not identically zero, then writing it as a power series, we have 
%     \begin{align*}
%         f(w)=(w-z)^ng(w)
%     \end{align*}
%     where $g(w)=a_n+a_{n+1}(w-z)+a_{n+2}(w-z)^2+\dots$, for some $a_n\neq 0$. We know $g(z)=a_n\neq 0$, hence there exists a neighborhood $V$ around $z$ such that $g(w)\neq 0$ whenever $w\in V$. Hence, if $w\in V$, then $f(w)=0$ iff $w=z$. This is a contradiction to $z_n\to z$ and $f(z_n)=0$ for all $n$.
% \end{proof}

% \begin{warn}
%     This shows that a nonzero holomorphic function has isolated zeros.
% \end{warn}

\begin{prob}[Spring 2018, Q4]
    Let $f$ be an entire function, suppose that for each $z_0$, the power series expansion 
    \begin{equation*}
        f(z)=\sum_{n=0}^\infty a_n(z-z_0)^n
    \end{equation*}
    has at least one coefficient $a_n=0$. Show that $f$ must be a polynomial.
\end{prob}
\begin{proof}
    It suffices to show that for some $k\in\mathbb{N}$, $f^{(k)}\equiv 0$ because
    \begin{equation*}
        a_n=\frac{1}{n!}f^{(n)}(z_0)
    \end{equation*}
    We first observe that 
    \begin{equation*}
        \bigcup_{n\in\mathbb{N}}\{z: f^{(n)}=0\}=\C
    \end{equation*}
    Then there must exist some $k\in\mathbb{N}$ such that 
    \begin{equation*}
        \{z: f^{(k)}=0\} \text{ is uncountable}
    \end{equation*}
    Thus $f^{(k)}$ has uncountably many zeros, which implies $f\equiv 0$. 
    \begin{lem}
        Let $f$ be a nonzero holomorphic function, then $f$ cannot have uncountably many zeros.
    \end{lem}
    \begin{proof}
        Let $Z=\{z\in\C: f(z)=0\}$, we know that $Z$ contains no limit point. We can tile the complex plane using unit squares and suppose $Z$ is uncountable, the intersection of $Z$ with one of the tiles needs to be uncountable. This is a contradiction because this intersection is compact and discrete, hence must be finite.
    \end{proof}
\end{proof}

\begin{prob}[Fall 2019, Q2]
    Let $f$ be an entire function such that 
    \begin{equation*}
        \sup_{|z|>R}|f(z)|\leq AR^k+B
    \end{equation*}
    for positive constants $A,B$ and all $R$ large. Then show $f$ is a polynomial with degree at most $k$.
\end{prob}
\begin{proof}
    This problem is a direct application of Cauchy's inequality. Writing $f(z)$ as a power series around $0$, we have 
    \begin{equation*}
        f(z)=\sum_{n=0}^\infty a_nz^n
    \end{equation*}
    where 
    \begin{equation*}
        a_n=\frac{1}{n!}f^{(n)}(0)
    \end{equation*}
    It suffices to show that for any $m>k$, we have 
    \begin{equation*}
        f^{(m)}(0)= 0
    \end{equation*}
    Using Cauchy's inequality, we get 
    \begin{align*}
        |f^{(m)}(0)|&\leq\frac{m!\sup_{|z|=R}|f(z)|}{R^m}\\
        &\leq \frac{m!(AR^k+B)}{R^m}\to 0
    \end{align*}
    as $R\to\infty$ because $m>k$. Hence we are done.
\end{proof}


\begin{prob}[Fall 2023, Q1]
    \begin{enumerate}
        \item[(a)] Show that if an entire function $f$ satisfies
        \[
            |f(z)| \leq \sqrt{|z|}
        \]
        for all $z \in \mathbb{C}$, then $f(z)$ is a constant function.
    
        \item[(b)] Why isn't the function $z^{1/2}$ a counterexample for Part (a)?
    \end{enumerate}
\end{prob}
\begin{proof}
    \begin{enumerate}
        \item[(a)] We write $f$ as the Taylor series about $0$, 
        \begin{equation*}
            f(z)=\sum_{n=0}^\infty a_nz^n
        \end{equation*}
        where 
        \begin{equation*}
            a_n=\frac{1}{n!}f^{(n)}(0)
        \end{equation*}
        Using Cauchy's inequality, for any $R$ large, if $n\geq 1$,
        \begin{align*}
            |f^{(n)}(0)|\leq\frac{n!\sup_{|z|=R}|f(z)|}{R^n}\to 0
        \end{align*}
        as $R\to\infty$. Thus $f(z)=a_0$, for some $a_0\in\C$.
        \item[(b)] The function is not entire.
    \end{enumerate}
\end{proof}


\begin{prob}[Fall 2018, Q4]
    Suppose $f,g$ are both entire functions such that $|f(z)|\leq |g(z)|$ for all $z\in\C$. Show that there exists $c\in\C$ such that $f=cg$.
\end{prob}
\begin{proof}
    We will use Liouville's theorem. It suffices to show that $h=f/g$ continuously extends to an entire function. Suppose $g(z_0)=0$ for some $z_0$, then we know $f(z_0)=0$, hence writing 
    \begin{equation*}
        f(z)=(z-z_0)^ns(z), \quad g(z)=(z-z_0)^m t(z)
    \end{equation*}
    where $s(z_0)\neq 0, t(z_0)\neq 0$. We must have $n\geq m$, suppose if not, i.e., $m>n$, then 
    \begin{equation*}
        |s(z)|\leq|z-z_0|^{m-n}|t(z)|
    \end{equation*}
    which is a contradiction by letting $z=z_0$. Thus $z_0$ is a removable singularity for $h=f/g$. This implies $h$ can be extended to an entire function, as desired.
\end{proof}


\begin{prob}
    
\end{prob}

\chapter{Residue formula}
\begin{prob}
    Compute 
    \begin{equation*}
        \int_{0}^\infty\frac{1}{1+x^2}dx
    \end{equation*}
\end{prob}


\begin{prob}
    Compute the following integral: fix $0<a<1$, 
    \begin{equation*}
        \int_{-\infty}^\infty\frac{e^{ax}}{1+e^x}dx
    \end{equation*}
\end{prob}
% \begin{proof}
%     Use rectangles and the residue at $\pi i$.
% \end{proof}





\begin{prob} 
    Assume that $f:\mathbb{C}\to\mathbb{C}$ is an entire function, not identically equal to $0$, and let \[ Z=\{z\in\mathbb{C}:f(z)=0\}. \] Prove that if $Z$ is unbounded, then $f$ has an essential singularity at $\infty$.
\end{prob}





\chapter{Spring 2026}
\begin{prob}
Let $f$ be a holomorphic function on $\{z:|z|<1\}$ that extends to a continuous
function on $\{z:|z|\leq 1\}$ and maps the closed unit disk to itself.
Furthermore, suppose
\[
f(0)=f'(0)=\cdots=f^{(n-1)}(0)=0.
\]
Show that:
\begin{enumerate}
    \item[(1)]
    \[
    |f(z)|\leq |z|^n \qquad \text{for all } |z|\leq 1.
    \]

    \item[(2)] If there is some $P\in D(0,1)\setminus\{0\}$ such that
    \[
    |f(P)|=|P|^n
    \]
    or if
    \[
    |f^{(n)}(0)|=n!,
    \]
    then
    \[
    f(z)=e^{i\theta}z^n.
    \]
\end{enumerate}
\end{prob}

\begin{prob}[Also Fall 2023, Q2]
Let $-1<a<0$. Compute
\[
\int_0^\infty \frac{x^a}{1+x}\,dx.
\]
\end{prob}

\begin{prob}
Two parts.

\begin{enumerate}
    \item[(1)] Let
    \[
    f:U=\{z:0<|z|<1\}\longrightarrow\mathbb{C}
    \]
    be holomorphic and nonconstant. If there exists a sequence
    $\{z_n\}\subset U$ that converges to $0$ and satisfies
    \[
    f(z_n)=0,
    \]
    then prove that $0$ is an essential singularity of $f$.

    \item[(2)] Show that if $f$ is an entire function that satisfies
    \[
    \lim_{|z|\to\infty}f(z)=\infty,
    \]
    then $f$ is a polynomial.
\end{enumerate}
\end{prob}
% \begin{proof}
%     Method 1: let $x=e^u$.
 
%     Method 2: use a keyhole contour.
%  \end{proof}

\begin{prob}
Describe all injective holomorphic maps that map
\[
\mathbb{C}\cup\{\infty\}
\quad\text{to}\quad
\mathbb{C}\cup\{\infty\}
\]
and also map
\[
\mathbb{R}\cup\{\infty\}
\quad\text{to}\quad
\mathbb{R}\cup\{\infty\}.
\]
\end{prob}



\end{document}